\subsection{Uniform structure and completion of a topological vector space}

Since the topology of the topological vector space E is compatible with the additive group structure on E, it defines a uniform structure on E (GT, III, § 3); when we speak of the uniform structure of a topological vector space we always mean this structure unless the contrary is expressly stated. Every continuous linear mapping of a topological vector space E in a topological vector space F is uniformly continuous (GT, III, § 3.1, prop. 3); every mapping of E in itself of the form \(x \mapsto ax + b\) is uniformly continuous. An equicontinuous set of linear mappings of E in F is uniformly equicontinuous (GT, X, § 2.2, prop. 5).

Remarks. --- 1) If B is a precompact set of K, then for every neighbourhood V of 0 in E, there is a neighbourhood U of 0 in E such that BU ⊂ V. For, if W is a neighbourhood of 0 in E such that W + W ⊂ V; then from (EVT \(_{III}\) ) there is a neighbourhood T \(_{0}\) of 0 in K and a neighbourhood U \(_{0}\) of 0 in E such that T \(_{0}\) U \(_{0}\) ⊂ W. As B is precompact, there are finitely many points \(\lambda_{i} \in \mathbf{B}\) ( \(1 \leqslant i \leqslant n\) ) such that the \(\lambda_{i} + T_{0}\) cover B; from (EVT \(_{II}\) ) it follows that there is a neighbourhood U ⊂ U \(_{0}\) of 0 in E, such that \(\lambda_{i}\) U ⊂ W for all i; clearly U has the required properties. In a similar manner (using (EVT \(_{I}\) ) instead of (EVT \(_{II}\) )) it can be shown that if H is a precompact set of E, then for every neighbourhood V of 0 in E, there exists a neighbourhood T of 0 in K such that TH ⊂ V.

\begin{enumerate}
\def\labelenumi{\arabic{enumi})}
\setcounter{enumi}{1}
\tightlist
\item
  From 1) it follows that, if B is a precompact set of K and H is a precompact set of E, then the mapping \((\lambda, x) \mapsto \lambda x\) restricted to \(\mathbf{B} \times \mathbf{H}\) is uniformly continuous. For, if V is a neighbourhood of 0 in E then there are neighbourhoods T of 0 in K, and U of 0 in E such that \(\mathrm{TH} + \mathrm{BU} \subset \mathrm{V}\) . Since we can write \(\lambda x - \lambda' x' = (\lambda - \lambda') x + \lambda'(x - x')\) , we see that for \(\lambda, \lambda'\) in B, \(x, x'\) in H, \(\lambda - \lambda' \in \mathrm{T}\) and \(x - x' \in \mathrm{U}\) , we have \(\lambda x - \lambda' x' \in \mathrm{V}\) , which proves our assertion.
\end{enumerate}

A topological vector space is called complete if, considering its uniform structure, it is a complete uniform space.

DEFINITION 2. --- A complete normed space on a non-discrete valued division ring is called a Banach space.

Examples. --- If K is a non-discrete valued division ring then the space \(\mathcal{B}(\mathrm{I};\mathrm{K})\) (I, p.~4, Example) is complete (GT, X, § 3.1, cor. 1). This is also true for the space \(\ell_{\mathbf{K}}^{1}(\mathrm{I})\) (I, p.~4, Example) with the norm \(\| x\| _1 = \sum_{\iota \in \mathrm{I}}|\xi_\iota |:\) for, if \(x_{n}\) is a Cauchy sequence in this space and \(x_{n} = (\xi_{n\iota})_{\iota \in \mathrm{I}}\) , then for all \(\iota \in \mathrm{I}\)

\[
\left| \xi_ {m\iota} - \xi_ {n\iota} \right| \leqslant \left\| x _ {m} - x _ {n} \right\| _ {1};
\]

thus, for each \(\iota \in I\) , the sequence \((\xi_{n\iota})_{n\geqslant 1}\) converges to a limit \(\xi_{\iota}\) in K. Further, for each finite subset J of I

\[
\sum_ {\iota \in J} | \xi_ {m\iota} - \xi_ {n\iota} | \leqslant \| x _ {m} - x _ {n} \| _ {1};
\]

and it follows immediately that there exists a constant \(a > 0\) , independent of J, \(m, n\) such that \(\sum_{\iota \in \mathbf{J}} |\xi_{m\iota} - \xi_{n\iota}| \leqslant a\) . Letting \(m\) tend to \(+\infty\) , we deduce \(\sum_{\iota \in \mathbf{J}} |\xi_{\iota} - \xi_{n\iota}| \leqslant a\) from which \(\sum_{\iota \in \mathbf{I}} |\xi_{\iota}| \leqslant a + \| x_n\|_1\) , which shows that \(z = (\xi_{\iota})_{\iota \in \mathbf{I}}\) belongs to \(\ell_{\mathbf{K}}^1 (\mathbf{I})\) ; further, for all \(\varepsilon > 0\) , there exists \(n_0\) such that for \(n \geqslant n_0\) and for every finite set J of I, we have \(\sum_{\iota \in \mathbf{J}} |\xi_{\iota} - \xi_{n\iota}| \leqslant \varepsilon\) ; passing to the limit with respect to the directed set of finite subsets of I, we see that \(\| z - x_n\| _1\leqslant \varepsilon\) for \(n\geqslant n_0\) , which shows that \(z\) is the limit of the sequence \((x_{n})\) in the normed space \(\ell_{\mathbf{K}}^{1}(\mathrm{I})\)

Let K be a Hausdorff topological division ring, E a topological vector space on K and suppose that the completed ring \(\hat{K}\) is a division ring (this is so when K is a valued division ring, GT, IX, § 3.3) then the Hausdorff completion \(\hat{E}\) of E carries the structure of a complete topological vector space on \(\hat{K}\) (GT, III, § 6.5); we say that \(\hat{E}\) , with this structure, is the Hausdorff completion of the topological vector space E, or simply the completion of E when E is Hausdorff.

